%%%%%%%%%%%%%%%%%%%%%%
% Marco teorico y antecedentes
%%%%%%%%%%%%%%%%%%%%%%

Este capítulo sienta las bases conceptuales que sostienen el resto del documento. En primer lugar, se delimitan los conceptos centrales —deriva arquitectónica, deuda técnica arquitectónica y \textit{smells}— para, a continuación, presentar las \textit{Fitness Functions} como instrumento de verificación, el análisis estático y los grafos de dependencia como insumo estructural, las métricas de acoplamiento que permiten cuantificarlos y la noción de gobernanza continua como marco de aplicación. Finalmente, se revisan los trabajos previos más relevantes y se contrastan mediante una tabla comparativa, con el fin de evidenciar el vacío que este proyecto pretende ocupar.

\section{Marco Teórico}

\subsection{Arquitectura de software y atributos de calidad}
Antes de abordar la deriva, conviene precisar qué se entiende por arquitectura de software y por qué su integridad importa. La arquitectura es el conjunto de estructuras principales de un sistema, junto con las relaciones que las conectan y los principios que gobiernan su diseño y evolución \citep{Bass2012}. No es un producto tangible, sino una propiedad emergente que se materializa en las dependencias, las capas, los módulos y los contratos que articulan el código. Su papel es doble: por un lado, sirve de vehículo para alcanzar los atributos de calidad deseados; por otro, actúa como marco de comunicación entre quienes construyen y quienes mantienen el sistema \citep{PerryWolf1992}.

La relevancia de la arquitectura para la mantenibilidad ha sido documentada de forma consistente. Cuando la estructura se degrada, los atributos de calidad no decaen todos a la vez ni de manera uniforme: la mantenibilidad y la evolucionabilidad suelen ser los primeros en resentirse, precisamente porque dependen de que las dependencias sean predecibles y de que los cambios puedan acotarse. Un sistema cuya arquitectura se ha erosionado obliga a sus desarrolladores a razonar sobre un territorio que ya no corresponde a ningún mapa compartido, lo que eleva el costo cognitivo y material de cada modificación \citep{Klein2022}. Esta relación entre estructura y costo de cambio es la razón de fondo por la que la vigilancia arquitectónica no puede considerarse un lujo.

\subsection{Deriva y erosión arquitectónica}
La deriva arquitectónica describe el distanciamiento gradual entre la arquitectura concebida para un sistema y la que efectivamente se materializa en su implementación \citep{PerryWolf1992, Li2022}. Se distingue de la erosión arquitectónica, término que suele reservarse para el deterioro estructural ya consumado y observable: la deriva es el proceso, la erosión es su huella acumulada. La distinción y sus múltiples denominaciones fueron sistematizadas por \citet{DeSilva2012}, quienes al revisar las estrategias disponibles para controlar la erosión las agruparon en tres familias —minimizarla mediante disciplina de proceso, prevenirla mediante restricciones impuestas por construcción y repararla una vez detectada— y señalaron que las aproximaciones basadas en verificación automática de conformidad se sitúan entre la primera y la segunda. Ambos conceptos comparten una misma raíz: las decisiones de implementación que, tomadas de forma incremental y bajo presión de entrega, van apartando al código de las restricciones estructurales que deberían gobernarlo \citep{Anthony2024}. La literatura coincide en señalar que este fenómeno no es un accidente puntual, sino una consecuencia casi inevitable de la evolución de los sistemas, lo que convierte su vigilancia en una tarea permanente y no en una actividad esporádica \citep{Klein2022}. Esta lectura es coherente con los trabajos clásicos sobre mantenimiento de la consistencia arquitectónica, que ya advertían del costo de permitir que la implementación y el diseño se separen \citep{Rosik}, y con las leyes de evolución del software, que prevén la degradación estructural como desenlace natural de todo sistema que crece \citep{Lehman1980}. De ahí que la deriva no deba tratarse como un fallo excepcional del proceso, sino como un riesgo estructural que reclama mecanismos de vigilancia integrados en el propio ciclo de desarrollo. El estudio de mapeo sistemático de \citet{Li2022}, que reunió 73 estudios primarios publicados entre 2006 y 2019, ofrece además una orientación concreta sobre qué vigilar: de las cuatro familias de síntomas de erosión que identifica ---estructurales, de violación, de calidad y de evolución---, las más frecuentes son las estructurales y las de violación, y entre sus ejemplos característicos figuran las dependencias cíclicas y la violación del patrón de capas. Son precisamente dos de las tres dimensiones que el anteproyecto fijó para este trabajo.

\subsection{Deuda técnica arquitectónica y \textit{smells}}
El concepto de deuda técnica, acuñado originalmente para describir el costo de las decisiones de diseño postergadas \citep{Cunningham1992}, trasladado al plano estructural, da lugar a la deuda técnica arquitectónica: el costo diferido que una organización asume al postergar la corrección de las desviaciones que degradan la arquitectura. Los \textit{smells} arquitectónicos —ciclos de dependencia, componentes excesivamente acoplados, violaciones de capas, concentraciones de responsabilidad— son los síntomas observables de esa deuda \citep{SasAvgeriou, ArcelliFontana2016}. Su costo no es solo teórico: en un estudio con desarrolladores de la industria, replicado y validado con un segundo conjunto de datos independiente, \citet{Besker2019} estimaron que la deuda técnica consume en promedio el $23\%$ del tiempo de desarrollo. La evidencia empírica ha establecido una relación negativa entre la presencia de estos \textit{smells} y la mantenibilidad de los sistemas, lo que refuerza la necesidad de instrumentos que los detecten de forma objetiva \citep{Jolak2025, Sas2019}. Para las PyMES, cuya capacidad de reacción ante la deuda acumulada es limitada, el reto no es solo detectar los síntomas, sino hacerlo antes de que su corrección resulte prohibitiva \citep{Hamed2023}.

\subsection{\textit{Fitness Functions}}
Una \textit{Fitness Function} es una verificación objetiva y automatizable que expresa, en términos cuantificables, una propiedad que el sistema debe preservar. El término, tomado de la biología evolutiva, fue adoptado por la ingeniería de software para describir pruebas que no verifican una funcionalidad, sino una característica estructural o de calidad de la arquitectura \citep{PaulWang2019}. La formulación de referencia es la de \citet{FordParsonsKua2017}, quienes las sitúan en el centro de las llamadas arquitecturas evolutivas y proponen clasificarlas según varios ejes: por su alcance, distinguiendo las \textit{atómicas} —que evalúan una única dimensión, como una regla de dependencia entre módulos— de las \textit{holísticas}, que evalúan la interacción de varias; y por su cadencia, separando las que se ejecutan de forma continua en cada integración de las que se activan de manera puntual o desencadenada por un evento. Las tres reglas que este trabajo implementa son, en esa taxonomía, funciones atómicas de ejecución desencadenada, y la Tasa de Conformidad que las agrega constituye una lectura holística construida sobre ellas. A diferencia de las métricas convencionales, que se limitan a describir el estado del sistema, las \textit{Fitness Functions} incorporan un criterio de juicio: producen un veredicto de cumplimiento o incumplimiento sobre una regla declarada. Esta propiedad las hace idóneas para la gobernanza arquitectónica, pues permiten codificar las decisiones de diseño como comprobaciones ejecutables que se revalúan en cada iteración del desarrollo. La segunda edición de la obra de referencia amplía precisamente ese uso, al tratar las \textit{Fitness Functions} como el instrumento de una gobernanza arquitectónica automatizada \citep{Ford2022}. Herramientas como \textit{ArchUnit} han popularizado esta idea en el ecosistema Java, al permitir expresar reglas arquitectónicas como pruebas unitarias \citep{ArchUnit}. En la práctica industrial se distingue, además, entre reglas que bloquean una integración de inmediato y aquellas que solo registran la desviación para su revisión posterior, distinción que la literatura técnica de divulgación denomina \textit{guardrails} arquitectónicos \citep{StefanoRossiniMedium}. El énfasis en convertir decisiones de diseño en comprobaciones conecta las \textit{Fitness Functions} con la documentación explícita de las decisiones arquitectónicas, cuya justificación ---hacer transparentes la arquitectura y su razón de ser para todos los interesados--- formularon \citet{TyreeAkerman2005}, y que la práctica actual registra en \textit{Architecture Decision Records} (ADR) \citep{ADR}. Este trabajo registra así sus propias decisiones (anexo~\ref{anexo:decisiones}).

\subsection{Análisis estático y grafos de dependencia}
El análisis estático consiste en examinar el código fuente —o sus artefactos compilados— sin ejecutar el programa, con el fin de extraer propiedades estructurales. Cuando el resultado de ese examen se organiza como un grafo dirigido, en el que los nodos representan componentes (paquetes, clases o módulos) y las aristas representan dependencias entre ellos, se obtiene una representación compacta y verificable de la arquitectura real del sistema \citep{Manglaras2026}. Este enfoque es particularmente valioso en el contexto de este trabajo porque permite prescindir de la documentación arquitectónica: la estructura se infiere directamente del código, que es la única fuente de verdad siempre disponible. La elección entre analizar el código fuente o el \textit{bytecode} compilado responde a un compromiso práctico: el primero ofrece mayor riqueza semántica, mientras que el segundo simplifica la resolución de dependencias al operar sobre una representación ya consolidada. En el ecosistema Java, herramientas como \textit{ArchUnit} ejemplifican este segundo enfoque, al importar clases compiladas y construir sobre ellas un modelo navegable de dependencias \citep{ArchUnit}, al igual que \textit{JDepend}, que calcula las métricas de paquetes de Martin directamente sobre el \textit{bytecode} \citep{JDepend}, y que \textit{Arcan}, que detecta \textit{smells} arquitectónicos sobre grafos extraídos de proyectos Java compilados \citep{ArcelliFontana2016}. Los modelos de reflexión de \citet{Murphy2001} y los mecanismos de sincronización arquitectónica \citep{ArchSyncJAndres} constituyen antecedentes directos de esta manera de contrastar el código con las expectativas estructurales.

Representar la estructura como un grafo dirigido no es una comodidad de notación: habilita el uso de algoritmos con propiedades conocidas. La detección de ciclos, en particular, se reduce a hallar los componentes fuertemente conexos del grafo, problema que \citet{Tarjan1972} resolvió con un algoritmo de una sola pasada en profundidad y complejidad lineal en el número de nodos y aristas. Un componente fuertemente conexo de tamaño mayor que uno es, por definición, un conjunto de paquetes mutuamente alcanzables, es decir, un ciclo. Esta correspondencia entre un concepto arquitectónico y una construcción de la teoría de grafos es la que permite que la verificación sea exacta y no heurística.

\subsection{Métricas de acoplamiento}
Algunas propiedades estructurales no se expresan como una prohibición binaria sino como una cuestión de grado, y para ellas hace falta una medida. \citet{Martin2003} introdujo dos métricas que se han vuelto estándar en el análisis de dependencias entre paquetes: el \textbf{acoplamiento eferente} ($Ce$), que cuenta cuántos paquetes distintos conoce un paquete dado, y el \textbf{acoplamiento aferente} ($Ca$), que cuenta cuántos dependen de él. La primera mide la exposición del paquete a los cambios ajenos; la segunda, su responsabilidad sobre los cambios propios.

Un valor elevado de $Ce$ señala un paquete que ha ido acumulando conocimiento de todo el sistema, y es el indicio habitual de lo que la literatura denomina componente dios: un módulo del que todo depende y que depende de todo, cuya modificación repercute en cascada y cuyo reemplazo se vuelve inviable \citep{SasAvgeriou}. Conviene subrayar, no obstante, que la métrica es puramente topológica y no distingue el papel que el paquete desempeña: una raíz de composición o una fachada presentan por diseño un $Ce$ alto sin que ello constituya una deriva. De ahí que el umbral a partir del cual un valor se considera excesivo no sea universal, sino un parámetro que debe calibrarse según el sistema y el rol de cada componente. Para esa calibración existe un procedimiento establecido: derivar el umbral de la distribución que la métrica presenta en sistemas de referencia y fijarlo en un percentil alto, en lugar de elegirlo a juicio \citep{LanzaMarinescu2006}. \citet{Alves2010} lo sistematizaron sobre cien sistemas y propusieron los percentiles 70, 80 y 90 como fronteras de riesgo moderado, alto y muy alto. Es el procedimiento que se aplica en la sección~\ref{sec:umbral} para fijar el umbral de concentración.

\subsection{Conformidad arquitectónica y gobernanza continua}
La conformidad arquitectónica es el grado en que la implementación de un sistema se adhiere a las reglas estructurales que se han definido para él \citep{KnodelPopescu2007, Klein2022}. Su verificación no es un fin en sí mismo, sino el insumo de una práctica más amplia: la gobernanza continua, entendida como la vigilancia sistemática y automatizada de las propiedades arquitectónicas a lo largo de todo el ciclo de vida del producto. En organizaciones pequeñas, donde los marcos de control deben ser livianos y proporcionados a los recursos disponibles, la gobernanza continua solo es viable si se apoya en instrumentos de bajo costo y de interpretación sencilla \citep{LAIG2025}. Esta constatación orienta el diseño del mecanismo que se describe en los capítulos siguientes: un instrumento que agregue las verificaciones individuales en una señal global comprensible para un equipo de desarrollo sin especialistas dedicados. A la vigilancia automatizada se suma, como complemento, la documentación de decisiones arquitectónicas. Los \textit{ADR} proporcionan el registro deliberado de por qué se adoptó una estructura y no otra, y constituyen el insumo natural contra el cual contrastar las desviaciones detectadas por el mecanismo \citep{ADR}. De esta manera, la gobernanza continua no se reduce a un acto de verificación, sino que articula un ciclo de documentación, monitoreo y corrección.

\section{Estado del Arte}
La detección de la deriva arquitectónica ha sido abordada desde múltiples frentes, que conviene examinar para situar la contribución de este proyecto. Un primer grupo de trabajos se concentra en la verificación de conformidad mediante modelos explícitos. Los modelos de reflexión de \citet{Murphy2001} comparan una representación de alto nivel con la estructura extraída del código, pero presuponen la existencia y el mantenimiento de ese modelo de alto nivel, un requisito que las PyMES rara vez satisfacen. \citet{KnodelPopescu2007} sistematizaron esta familia al comparar los tres enfoques disponibles para la verificación estática de conformidad —modelos de reflexión, matrices de dependencias y lenguajes de reglas— y concluyeron que ninguno domina a los demás en todas las dimensiones: los modelos de reflexión ofrecen mayor poder expresivo a cambio de exigir una arquitectura de referencia, mientras que los lenguajes de reglas resultan más económicos de adoptar pero cubren un espectro menor de restricciones. Esa constatación es directamente pertinente aquí, pues el presente trabajo se sitúa en la tercera familia y asume conscientemente ese compromiso. En la misma línea, \citet{ArchSyncJAndres} exploran la sincronización entre la arquitectura documentada y la implementación, y \citet{SanMartin}, en un trabajo aún en forma de preprint, estudian la verificación de conformidad en sistemas autoadaptativos mediante un lenguaje específico para declarar la arquitectura planificada; ambos asumen igualmente una especificación arquitectónica de referencia.

Un segundo grupo aborda la deriva a través de sus síntomas. \citet{ArcelliFontana2016} detectan automáticamente \textit{smells} de inestabilidad ---dependencias cíclicas, concentradoras (\textit{hub-like}) e inestables--- sobre el grafo de dependencias de proyectos compilados, \citet{Sas2019} investigan su evolución, y \citet{SasAvgeriou} proponen un índice de deuda técnica arquitectónica basado en aprendizaje de máquina aplicado a la detección de \textit{smells}. Estos trabajos aportan evidencia valiosa sobre el diagnóstico, pero se orientan a la caracterización del deterioro más que a ofrecer un veredicto de conformidad accionable para un equipo de desarrollo cotidiano.

Un tercer grupo lo constituyen las propuestas centradas en las \textit{Fitness Functions} como práctica. \citet{PaulWang2019} y \citet{Ford2022} formalizan la idea de verificar propiedades arquitectónicas mediante comprobaciones automatizadas, y \textit{ArchUnit} la materializa en el ecosistema Java \citep{ArchUnit}. Estos trabajos definen el \textit{cómo} verificar, pero no proponen un mecanismo que agregue los resultados en una clasificación global del estado del sistema. Por último, trabajos recientes como \citet{Manglaras2026} reconstruyen arquitecturas de sistemas distribuidos a partir del código, lo que confirma el interés creciente por extraer estructura sin depender de documentación.

Un cuarto grupo, más reciente, agrupa las herramientas y guías orientadas a hacer práctica la verificación en el día a día de los equipos. Los \textit{guardrails} arquitectónicos, difundidos como complemento de las \textit{Fitness Functions} \citep{StefanoRossiniMedium}, y las soluciones de verificación por convención para distintos ecosistemas —como \textit{NetArchTest} para .NET \citep{SugumarMedium} o \textit{ArchUnit} para Java \citep{ArchUnitCodecentric}— evidencian que la industria ya demanda instrumentos ligeros de control estructural. Estas fuentes se citan como testimonio de la práctica, no como respaldo de las decisiones del trabajo, que descansan en la literatura arbitrada. Sin embargo, estas herramientas comparten una limitación: exigen que alguien escriba y mantenga las reglas, y no ofrecen por sí solas una lectura agregada del estado global del sistema.

La tabla~\ref{tab:estadoArte} sintetiza esta comparación en torno a cuatro criterios: el enfoque adoptado, si la propuesta detecta derivas de forma automatizada, si prescinde de documentación arquitectónica externa y si resulta asequible para una PyME sin especialistas dedicados.

\begin{table}[H]
\centering
\caption{Comparación de los trabajos relacionados con la propuesta. Leyenda: $\checkmark$ lo cumple, $\sim$ lo cumple parcialmente, $\times$ no lo cumple.}
\label{tab:estadoArte}
\footnotesize
\begin{tabularx}{\textwidth}{>{\raggedright\arraybackslash}p{3.7cm}>{\raggedright\arraybackslash}Xcccc}
\toprule
\textbf{Trabajo} & \textbf{Enfoque} &
\rotatebox[origin=l]{90}{\textbf{Automática}} &
\rotatebox[origin=l]{90}{\textbf{Sin doc.\ externa}} &
\rotatebox[origin=l]{90}{\textbf{Veredicto agregado}} &
\rotatebox[origin=l]{90}{\textbf{Asequible PyME}} \\
\midrule
\citet{Murphy2001} & Modelos de reflexión & $\times$ & $\times$ & $\times$ & $\times$ \\
\citet{KnodelPopescu2007} & Comparación de enfoques de conformidad & $\checkmark$ & $\times$ & $\times$ & $\sim$ \\
\citet{ArchSyncJAndres} & Sincronización arquitectónica & $\sim$ & $\times$ & $\times$ & $\times$ \\
\citet{SanMartin} & Verificación de conformidad (MAPE-K) & $\checkmark$ & $\times$ & $\times$ & $\times$ \\
\citet{ArcelliFontana2016} & Detección de \textit{smells} (\textit{Arcan}) & $\checkmark$ & $\checkmark$ & $\times$ & $\sim$ \\
\citet{Sas2019} & Evolución de \textit{smells} de inestabilidad & $\sim$ & $\checkmark$ & $\times$ & $\sim$ \\
\citet{SasAvgeriou} & Índice de deuda técnica con aprendizaje de máquina & $\sim$ & $\checkmark$ & $\checkmark$ & $\sim$ \\
\citet{FordParsonsKua2017, Ford2022}; \citet{PaulWang2019} & \textit{Fitness Functions} como práctica & $\checkmark$ & $\sim$ & $\times$ & $\sim$ \\
\citet{ArchUnit} & Reglas arquitectónicas como pruebas & $\checkmark$ & $\sim$ & $\times$ & $\sim$ \\
\citet{Manglaras2026} & Reconstrucción de arquitectura & $\checkmark$ & $\checkmark$ & $\times$ & $\sim$ \\
\midrule
\textbf{Esta propuesta} & \textit{Fitness Functions} sobre grafos, con tasa de conformidad y severidad & $\checkmark$ & $\checkmark$ & $\checkmark$ & $\checkmark$ \\
\bottomrule
\end{tabularx}
\end{table}

La columna «veredicto agregado» merece una precisión, porque es la que sostiene el argumento de la brecha. Solo \citet{SasAvgeriou} producen una cifra global, pero lo hacen mediante un modelo de aprendizaje de máquina cuyo resultado no es reconstruible por el equipo que lo recibe; el índice indica cuánta deuda hay, no qué reglas se incumplieron ni dónde. La agregación que aquí se propone persigue el objetivo contrario: una cifra deliberadamente simple, calculada como un cociente que cualquier miembro del equipo puede rehacer a mano a partir de los hallazgos que la componen.

El examen anterior revela una brecha consistente: los enfoques basados en modelos exigen documentación que las PyMES no mantienen; los enfoques basados en \textit{smells} describen el deterioro sin ofrecer un veredicto agregado; y las propuestas de \textit{Fitness Functions} aportan el mecanismo de verificación sin resolver la agregación en un diagnóstico clasificatorio. Ninguno de los trabajos revisados combina, en un solo instrumento ligero y sin dependencias de documentación externa, la extracción del grafo de dependencias, la evaluación mediante \textit{Fitness Functions} y la clasificación del estado de conformidad en niveles interpretables. Es precisamente ese espacio el que ocupa la presente investigación. Conviene precisar desde ahora en qué sentido la propuesta prescinde de la documentación: no de toda declaración de la intención ---la deriva solo puede medirse contra una arquitectura prevista---, sino de la documentación arquitectónica extensa que los enfoques basados en modelos presuponen, que se sustituye por un catálogo mínimo de capas y dependencias admitidas que vive junto al código y se verifica en cada ejecución.

\section{Resumen del capítulo}
Este capítulo definió los conceptos fundamentales que articulan el proyecto: deriva y erosión arquitectónica, deuda técnica arquitectónica y \textit{smells}, \textit{Fitness Functions}, análisis estático con grafos de dependencia, las métricas de acoplamiento que permiten cuantificar la concentración de responsabilidades y la gobernanza continua como marco de aplicación. Dos correspondencias resultan especialmente relevantes para los capítulos siguientes: la que vincula la noción de ciclo de dependencia con los componentes fuertemente conexos de un grafo, que permite verificarla de forma exacta \citep{Tarjan1972}, y la que vincula el componente dios con el acoplamiento eferente, que permite medirlo aunque exija calibrar su umbral \citep{Martin2003}. Sobre esa base, se revisaron los trabajos previos agrupándolos en cuatro corrientes —verificación basada en modelos, detección de \textit{smells}, \textit{Fitness Functions} como práctica y herramientas de verificación por convención— y se los contrastó mediante una tabla comparativa. La comparación evidenció la brecha que motiva la propuesta: la ausencia de un mecanismo que integre, de forma asequible y sin documentación externa, la detección y una clasificación agregada del estado de conformidad que siga siendo reconstruible por quien la recibe.